\subsection{Nullstellen, Linearfaktoren und Vielfachheit}
\label{subsec:real_functions_polynomial_zeros}
Die ausmultiplizierte Form zeigt die Koeffizienten. Eine Produktdarstellung
macht dagegen die Nullstellen sichtbar und wird später beim Kürzen von
Brüchen wichtig.
\begin{TheoremBox}{Nullstelle und Linearfaktor}{real_functions_factor_theorem}
  Für ein Polynom $p$ ist $x_0$ genau dann eine Nullstelle, wenn der
  \emph{Linearfaktor} \EnglishTerm{linear factor} $(x-x_0)$ abgespalten
  werden kann:
  \[
    p(x_0)=0
    \quad\Longleftrightarrow\quad
    p(x)=(x-x_0)q(x)
  \]
  mit einem Polynom $q$. Ist $p$ nicht das Nullpolynom und vom Grad $n$,
  so hat es höchstens $n$ verschiedene reelle Nullstellen.
\end{TheoremBox}

\begin{ExampleBox}{Eine bekannte Nullstelle nutzen}{real_functions_polynomial_division}
  Für $p(x)=x^3-x^2-x+1$ erkennt man $p(1)=0$.
  Die \emph{Polynomdivision} \EnglishTerm{polynomial division}
  durch $x-1$ verringert den Grad.
  \begin{ClozeArea}[34]{11}
    Nacheinander wird jeweils der höchste Term durch $x$ geteilt und das
    zugehörige Vielfache von $x-1$ abgezogen:
    \[
      \begin{array}{r@{\;}r@{\;}r@{\;}r@{\qquad}l}
        x^3&-x^2&-x&+1&:(x-1)=x^2-1\\
        -(x^3&-x^2)&&&\\ \cline{1-2}
        &&-x&+1&\\
        &&-(-x&+1)&\\ \cline{3-4}
        &&&0&\text{Rest}
      \end{array}
    \]
    Die binomische Formel liefert anschließend die vollständige Zerlegung:
    \[
      p(x)=(x-1)(x^2-1)=(x-1)^2(x+1).
    \]
  \end{ClozeArea}
  Die Nullstellen sind $\Cloze{4cm}{x=-1\ \text{und}\ x=1}$.
\end{ExampleBox}

\par\noindent\begin{minipage}{\linewidth}
\begin{DefinitionBox}{Vielfachheit einer Nullstelle}{real_functions_multiplicity}
  Eine Nullstelle $x_0$ eines Polynoms $p\ne0$ heißt
  \emph{$k$-fach} \EnglishTerm{of multiplicity $k$}, wenn
  \[
    p(x)=(x-x_0)^kq(x),\qquad q(x_0)\ne0.
  \]
  Der Faktor $(x-x_0)$ tritt also genau $k$-mal auf.
\end{DefinitionBox}
\begin{ExampleBox}[before skip=0.3\baselineskip]{Schneiden oder Berühren?}{real_functions_multiplicity_graph}
  Bei $p(x)=(x+1)(x-1)^2$ haben die beiden Nullstellen unterschiedliche
  Vielfachheiten.
  \par\smallskip
  \begin{minipage}[c]{0.46\linewidth}
    Bei $x=-1$ wechselt $x+1$ sein Vorzeichen: Der Graph
    $\Cloze{3cm}{\text{schneidet}}$ die $x$-Achse.

    Bei $x=1$ bleibt $(x-1)^2$ auf beiden Seiten positiv: Der Graph
    $\Cloze{3cm}{\text{berührt}}$ die $x$-Achse.
  \end{minipage}\hfill
  \begin{minipage}[c]{0.50\linewidth}
    \centering
    \CachedTikzFigure{content/resources/chapters/real_functions/tikz/functions_multiplicity.tex}
  \end{minipage}
\end{ExampleBox}
\end{minipage}

\begin{KeyPointBox}{Vielfachheit und Vorzeichenwechsel}{real_functions_multiplicity_sign}
  An einer reellen Nullstelle eines Polynoms gilt:
  \[
    \begin{array}{c@{\qquad}l}
      k\text{ ungerade}&\text{Vorzeichenwechsel: Achse wird geschnitten},\\
      k\text{ gerade}&\text{kein Vorzeichenwechsel: Achse wird berührt}.
    \end{array}
  \]
  Ein Polynom muss nicht so viele reelle Nullstellen wie sein Grad besitzen:
  $x^2+1$ hat beispielsweise keine reelle Nullstelle.
\end{KeyPointBox}
